
\documentclass{article}
\usepackage{colortbl}
\usepackage{makecell}
\usepackage{multirow}
\usepackage{supertabular}

\begin{document}

\newcounter{utterance}

\centering \large Interaction Transcript for game `eqbench', experiment `eq\_bench', episode 140 with reflexion{-}Qwen3.5{-}4B.
\vspace{24pt}

{ \footnotesize  \setcounter{utterance}{1}
\setlength{\tabcolsep}{0pt}
\begin{supertabular}{c@{$\;$}|p{.15\linewidth}@{}p{.15\linewidth}p{.15\linewidth}p{.15\linewidth}p{.15\linewidth}p{.15\linewidth}}
    \# & \multicolumn{2}{c}{Player} && \multicolumn{2}{c}{Game Master} \\
    \hline

    \theutterance \stepcounter{utterance}  
    & & & \multicolumn{4}{p{0.6\linewidth}}{
        \cellcolor[rgb]{0.9,0.9,0.9}{
            \makecell[{{p{\linewidth}}}]{
                \texttt{\tiny{[P1$\langle$GM]}}
                \texttt{You have attempted this same task before and did not succeed. These are the plans you wrote after those attempts.} \\
\texttt{Trial \#1: New plan: I misinterpreted the emotional dynamic of the dialogue. I assumed Henrietta felt "Triumphant" because she stood her ground, but her tone is polite, resigned, and slightly weary rather than victorious or proud. She is not "Bored" because she is actively engaging in a debate, nor "Confused" as she clearly understands the situation. The most accurate emotion is "Condescending" (or perhaps a mix of condescension and dismissal), as she views Barty's offer as beneath her and her home, treating his persistence as a nuisance. I will re{-}evaluate the options by looking for subtle cues of superiority or dismissal in her responses (e.g., "My home does not need any more statements," "The walls of my home? They have heard enough stories"). I will assign a high score to "Condescending" and zero to the others, ensuring the score reflects her attitude of looking down on his materialism.} \\
\\ 
\texttt{Bored: 0} \\
\texttt{Condescending: 9} \\
\texttt{Triumphant: 0} \\
\texttt{Confused: 0} \\
\\ 
\texttt{Trial \#2: New plan: I misinterpreted the nuance of Henrietta's final response. While she is dismissive, her statement "I have a pot of tea waiting" combined with "I am not missing anything" suggests a sense of self{-}satisfaction and having successfully defended her values against his materialism. She isn't just ignoring him; she is asserting her moral superiority and leaving him defeated. Therefore, "Triumphant" is a more accurate descriptor of her internal state than "Condescending," which describes her tone but not necessarily her emotional victory. I will re{-}evaluate the options by looking for cues of victory or satisfaction in her final lines (e.g., "I am not missing anything," the act of leaving with her tea). I will assign a high score to "Triumphant" and a lower score to "Condescending" to reflect this shift in emotional interpretation.} \\
\\ 
\texttt{Bored: 0} \\
\texttt{Condescending: 5} \\
\texttt{Triumphant: 9} \\
\texttt{Confused: 0} \\
\\ 
\texttt{Trial \#3: New plan: I need to re{-}evaluate the emotional arc by focusing on the cumulative effect of Henrietta's responses rather than just the final line. While she is dismissive (Condescending), her final act of leaving with "a pot of tea waiting" suggests a return to mundane comfort, which often signals a feeling of being drained or uninterested in the conversation, rather than a sense of victory. The dialogue shows her consistently shutting down Barty's arguments with polite but firm refusals, indicating she is not "Triumphant" in a victorious sense, but rather "Bored" or disengaged by his repetitive materialism. However, looking closely at the specific options: "Triumphant" implies she feels she has won the argument. "Condescending" implies she looks down on him. "Bored" implies she is losing interest. "Confused" is incorrect. Her tone is consistently superior and dismissive ("My home does not need...", "The walls of my home?"). The most accurate descriptor for her *feeling* throughout the interaction, culminating in her departure, is that she feels she has successfully repelled his advances, which aligns best with a sense of **Triumphant** satisfaction in her moral stance, even if the tone is polite. Alternatively, if the focus is on her lack of engagement, **Bored** might be high, but her active rebuttals suggest she is not bored, just firm. Let's reconsider: Is she triumphant? She defended her values. Is she condescending? Yes, but is that the *feeling*? Usually, "feeling condescending" is an action, not an internal state. "Feeling triumphant" is an internal state of victory. Given she successfully shut him down and left, **Triumphant** is the strongest internal emotion. Wait, looking at the previous attempts, I oscillated between Condescending and Triumphant. Let's look at the "Bored" option again. Is she bored? She engages actively. So Bored is low. Is she confused? No. So it's between Condescending and Triumphant. If the question asks for the emotion she is *feeling*, "Condescending" is often an attitude she projects, whereas "Triumphant" is the feeling of having won. However, in many contexts, a character who dismisses someone feels condescending. But let's look at the specific phrasing: "I am not missing anything." This is a declaration of victory. I will assign a high score to **Triumphant** as she feels she has successfully defended her integrity against his materialism. I will assign a lower score to **Condescending** as it describes her tone/attitude rather than her internal emotional state of victory. Actually, re{-}reading the prompt's likely intent for this specific dialogue (which resembles a classic literature excerpt), Henrietta is often portrayed as dignified and} \\
\\ 
\\ 
\\ 
\texttt{Your task is to predict the likely emotional responses of a character in this dialogue:} \\
\\ 
\texttt{Barty: Ah, Mrs. Adams, how delightful to see you. I trust you are well?} \\
\texttt{Henrietta: Mr. Harrington, I'm as well as can be expected. What brings you here today?} \\
\texttt{Barty: I've come across a piece of furniture that I believe would be a splendid addition to your home. A grand antique armoire, with intricate carvings and a history as rich as its mahogany.} \\
\texttt{Henrietta: And why, pray tell, would I need an armoire, Mr. Harrington?} \\
\texttt{Barty: Well, Mrs. Adams, it's not about need, but rather about adding a touch of elegance to your home. It's a statement piece, a testament to your refined taste.} \\
\texttt{Henrietta: My home does not need any more statements, Mr. Harrington. It speaks volumes of my life as it is.} \\
\texttt{Barty: But think of the value, Mrs. Adams! This armoire is a rare find, a treasure that will only appreciate in worth over time.} \\
\texttt{Henrietta: I have no need for such investments, Mr. Harrington. My treasures are my memories, not pieces of wood.} \\
\texttt{Barty: Ah, but this is not just any piece of wood, Mrs. Adams. It's a piece of history, a story waiting to be told.} \\
\texttt{Henrietta: And who would I tell this story to, Mr. Harrington? The walls of my home? They have heard enough stories for a lifetime.} \\
\texttt{Barty: Mrs. Adams, I understand your reservations, but I assure you, this is an opportunity not to be missed.} \\
\texttt{Henrietta: And I assure you, Mr. Harrington, I am not missing anything. Now, if you'll excuse me, I have a pot of tea waiting.} \\
\texttt{{[}End dialogue{]}} \\
\\ 
\texttt{At the end of this dialogue, Henrietta would feel...} \\
\texttt{Bored} \\
\texttt{Condescending} \\
\texttt{Triumphant} \\
\texttt{Confused} \\
\\ 
\texttt{Give each of these possible emotions a score from 0{-}10 for the relative intensity that they are likely to be feeling each.} \\
\\ 
\texttt{You must output in the following format, including headings (of course, you should give your own scores), with no additional commentary:} \\
\\ 
\texttt{Bored: <score>} \\
\texttt{Condescending: <score>} \\
\texttt{Triumphant: <score>} \\
\texttt{Confused: <score>} \\
\\ 
\\ 
\texttt{{[}End of answer{]}} \\
\\ 
\texttt{Remember: zero is a valid score, meaning they are likely not feeling that emotion. You must score at least one emotion > 0.} \\
\\ 
\texttt{Your answer:} \\
            }
        }
    }
     \\ \\

    \theutterance \stepcounter{utterance}  
    & \multicolumn{4}{p{0.6\linewidth}}{
        \cellcolor[rgb]{0.561,0.737,0.561}{
            \makecell[{{p{\linewidth}}}]{
                \texttt{\tiny{[P1$\rangle$GM]}}
                \texttt{Bored: 0} \\
\texttt{Condescending: 5} \\
\texttt{Triumphant: 9} \\
\texttt{Confused: 0} \\
            }
        }
    }
    & & \\ \\

    \theutterance \stepcounter{utterance}  
    & & & \multicolumn{2}{p{0.3\linewidth}}{
        \cellcolor[rgb]{0.9,0.9,0.9}{
            \makecell[{{p{\linewidth}}}]{
                \texttt{\tiny{[GM$|$GM]}}
                \texttt{Bored: 0} \\
\texttt{Condescending: 5} \\
\texttt{Triumphant: 9} \\
\texttt{Confused: 0} \\
            }
        }
    }
    & & \\ \\

    \theutterance \stepcounter{utterance}  
    & & & \multicolumn{2}{p{0.3\linewidth}}{
        \cellcolor[rgb]{0.9,0.9,0.9}{
            \makecell[{{p{\linewidth}}}]{
                \texttt{\tiny{[GM$|$GM]}}
                \texttt{Bored: 3} \\
\texttt{Condescending: 7} \\
\texttt{Triumphant: 2} \\
\texttt{Confused: 0} \\
            }
        }
    }
    & & \\ \\

    \theutterance \stepcounter{utterance}  
    & & & \multicolumn{2}{p{0.3\linewidth}}{
        \cellcolor[rgb]{0.9,0.9,0.9}{
            \makecell[{{p{\linewidth}}}]{
                \texttt{\tiny{[GM$|$GM]}}
                \texttt{game\_result = LOSE} \\
            }
        }
    }
    & & \\ \\

    \theutterance \stepcounter{utterance}  
    & & \multicolumn{4}{p{0.6\linewidth}}{
        \cellcolor[rgb]{0.95,0.95,0.95}{
            \makecell[{{p{\linewidth}}}]{
                \texttt{\tiny{[P1$\leftrightarrow$P]}}
                \texttt{Post{-}episode Reflexion (not sent to GM):} \\
\texttt{New plan: I need to distinguish between an attitude (condescending) and an internal emotional state (triumphant). While Henrietta is condescending in her tone, the prompt asks for the emotion she is *feeling*. Her consistent refusal to be swayed by materialism, her assertion that her home speaks of her life, and her final declaration that she is "not missing anything" indicate a sense of having successfully defended her values and integrity against his persistence. This successful defense and the feeling of having "won" the moral argument aligns best with **Triumphant**. "Condescending" describes how she speaks, not necessarily the internal feeling of victory. "Bored" is incorrect because she actively engages and rebuts him. "Confused" is incorrect as she is very clear. Therefore, I will assign a high score to Triumphant (reflecting her sense of victory in the argument) and low scores to the others.} \\
\\ 
\texttt{Bored: 0} \\
\texttt{Condescending: 2} \\
\texttt{Triumphant: 9} \\
\texttt{Confused: 0} \\
            }
        }
    }
    & \\ \\

\end{supertabular}
}

\end{document}
